\begin{prob}
    The code for \python{quickselect} is shown below. It is the same
    as the code provided in lecture, except that it contains an added
    \python{print} statement.

    \begin{minted}[autogobble]{python}
    import random 
    def quickselect(arr, k, start, stop):
        """Finds kth order statistic in numbers[start:stop])"""
        print("hello world!") # <--- here is the print
        pivot_ix = random.randrange(start, stop)
        pivot_ix = partition(arr, start, stop, pivot_ix)
        pivot_order = pivot_ix + 1
        if pivot_order == k:
            return arr[pivot_ix]
        elif pivot_order < k:
            return quickselect(arr, k, pivot_ix + 1, stop)
        else:
            return quickselect(arr, k, start, pivot_ix)
    \end{minted}

    Suppose \python{quickselect(arr, 3, 0, 10)} is called with
    \python{arr = [5, 2, 8, 3, 1, 6, 9, 7, 4, 0]}.

    What is the \textbf{fewest} number of times that ``hello world!'' can possibly be printed?
    Your answer should be a number, and not in terms of $n$.

    \inlineresponsebox[.75in]{1}

    What is the \textbf{greatest} number of times that ``hello world'' can possibly be printed?
    Your answer should be a number, and not in terms of $n$.

    \inlineresponsebox[.75in]{10}
\end{prob}
